\majorbreak
\section{Preserve trust and verify the migration}
\label{sec:validation}

A shorter proof is acceptable only if it proves the same proposition about
the same objects under the same trust policy. Static checks narrow the search
for good edits; they do not close that obligation. This section separates
the evidence obtained here from the gates still required before adoption.

\subsection{What the static audit actually establishes}

The comment/string-masked scan covers all 546 active Lean files. It found
no code-token hits for \code{sorry}, \code{admit}, \code{sorryAx}, explicit
\code{axiom} declarations, \code{native_decide}, \code{decide +native},
\code{bv_decide}, native-reflection symbols, or
\code{unsafe}/\code{extern}/\code{implemented_by}. It counted 2,163
\code{#print axioms} commands, zero \code{#guard_msgs} commands, and 525
\code{decide +kernel} occurrences.

These are \emph{lexical} findings. The scanner is not a Lean parser, and the
result says nothing conclusive about macro expansion, imported dependencies,
or the elaborated proof's transitive axiom set. External imports include
\code{Diophantine}, \code{BusyBeaver.ZFC}, \code{BoundedZFCConsistency},
\code{FirstOrder}, and Mathlib. Nor does the word \code{Native} in a module
name indicate use of Lean's native evaluator.

The finite circuit analysis independently checked list lengths, input and
fixed-leaf bounds, strictly earlier references, and every positive exponent
transition. An altered exponent and a forward/non-earlier edge were
rejected in disposable Python fixtures for both circuits. These are useful
data-consistency and checker-coverage tests, \textbf{not Lean proofs}.

The validation helpers passed 16 Python tests, including malformed/missing
axiom-report cases. The shell replay harness passed syntax checking and its
unavailable-toolchain branch returned the intended exit code 69. Its actual
Lean success branch was not run. Neither this evidence nor a successful
PDF build is a candidate elaboration result.

\subsection{Turn the existing audit output into a failing policy gate}
\label{sec:axiomgate}

The repository documents the intended standard-axiom policy and reproduction
commands at \docsrc{README.md}{908}{939}. Preserve it: an accepted endpoint
may depend on a subset of
\[
 \{\texttt{propext},\ \texttt{Classical.choice},\ \texttt{Quot.sound}\}.
\]
Reject every other axiom rather than checking only a few known bad names.
This catches custom assumptions, \code{sorryAx}, and native-evaluation
dependencies irrespective of their spelling in a particular Lean release.

\code{#print axioms} reports dependencies but does not itself reject a
violation. No failure-producing axiom policy was found in the inspected
committed locations: the pinned repository root has no \code{.github}
entry, and the inspected root scripts contain no such gate. This does not
establish that external or private CI is absent.

The supplied \pathcode{validation/check_axiom_log.py} rejects missing
expected targets, an empty target list, malformed/incomplete reports,
unexpected axioms, compiler-error text, and sorry warnings. It requires an
explicit endpoint list and accepts wrapped axiom lists and apostrophes in
names. Its successful status is deliberately \code{PASS_LOG_ONLY}: it
does not authenticate the log or establish a successful compiler exit.
Capture merged stdout/stderr and preserve Lean's exit code separately.

Apply the policy to both baseline and candidate endpoints, and to new helpers
on which they depend. Retain proof-producing finite methods, including the
project's \code{decide +kernel} convention. Replacing them with
\code{native_decide} to shorten source would change the trust boundary,
not achieve a like-for-like simplification.

\subsection{Consolidate pure audit wrappers but keep the regression proofs}

\textbf{Priority 2 for audit maintenance.} Of the 220 filename/path-defined
check or audit files, only 60 consist purely of import/open/namespace/end
commands and \code{#print axioms}. Those wrappers occupy 739 lines and
33,098 bytes and contain 530 print sites. A declarative manifest of fully
qualified endpoint names, required imports, test group, and axiom policy
could generate their scaffolding.

This does not eliminate 739 lines of mathematical proof: the names and all
530 checks must survive. The 217 already fully qualified printed targets
observed in the current source are all distinct, so no measured duplicated
target removal is claimed. Resolve unqualified names in their existing
namespace before centralizing them. Keep small audit shards for targeted
rebuilds rather than defaulting to one expensive umbrella test.

Do not delete the other checks as boilerplate. The full check/audit group
contains \textbf{835 examples and 85 theorem or lemma declarations}.
In particular:

\begin{itemize}
\item \src{NativeCyclicCarryAudit.lean}{34}{70} tests wraparound, a reset,
and genuine nonuniqueness; the last protects against an invalid uniqueness
assumption about cyclic carries.
\item \src{CurrentHighDigits2011Checks.lean}{10}{36} tests boundary cells
and origin subtraction. Its check at
\src{CurrentHighDigits2011Checks.lean}{48}{54} preserves the implication
from polynomial zero rather than assuming the desired mask property.
\item \src{Carrier65Checks.lean}{18}{49} fixes coordinate slots and
negative-tau behavior.
\end{itemize}

The selected \pathcode{EndpointAudit.lean} in this archive contains 35
additional target probes. It is uncompiled and not exhaustive; it does not
replace these existing behavioral checks.

\subsection{An incremental acceptance roadmap}
\label{sec:roadmap}

\textbf{Stage 0\quad Establish the baseline.}
Use a full checkout at the pinned commit, Lean 4.32.0, and Mathlib
\code{81a5d257c8e410db227a6665ed08f64fea08e997}. Build the original target
and run its existing audits, recording exact source/configuration hashes,
compiler exits, and raw logs. The presence of an old \code{.olean} file
is not evidence of a fresh trustworthy build. The focused archive lacks
the full dependency tree and cannot substitute for this baseline.

\textbf{Stage 1\quad Pilot the local reuse edits.}
Start with \code{half_base_bound_31}, finite-set support, and the
sum-of-squares proof. Keep every public declaration statement unchanged.
Then try the tuple helper and named accessors. Compile new generic modules
first, compatibility modules second, and every affected reverse dependency
afterward. This stage establishes the workflow before touching generators.

\textbf{Stage 2\quad Test the conservative circuit certificate.}
Compile the API-preserving power candidates against the unchanged literal
DAGs. Confirm all original public endpoints and recipe consumers, not just
the new private theorem. Compare the baseline and candidate axiom sets and
record actual certificate cost. Only then pilot the indexed facade on one
consumer. Update the authoritative generator if adopted; hand-editing only
generated output would lose the refactor on regeneration.

\textbf{Stage 3\quad Extract the shared semantic interfaces.}
Introduce the source/scalar-window certificate, refactor the two simple
digit clients, then the full window proof. Perform the base-module import
split only after that interface is stable. Next parameterize lower-layer
degree profiles, retaining concrete wrappers and all sharp constants.
Try the smaller coefficient-map cleanup before an AST migration.

\textbf{Stage 4\quad Accept or reject on evidence.}
Build the explicit main library target, the full affected import closure,
and all 220 existing check/audit files. Compare exported types, preserved
definitions or semantic bridges, axiom sets, boundary/negative tests, and
generator reproducibility. Keep the old names as aliases until consumers
pass. Reject a shorter edit if it hides a changed assumption, worsens
maintainability, or creates unacceptable replay cost.

\subsection{The invariants that a line-count target must not erase}

\begin{enumerate}
\item \textbf{Statement and definition preservation.}
Keep fully qualified public names, universe parameters, implicit/explicit
binders, typeclasses, hypotheses, and conclusions. Identical theorem text
is insufficient if referenced definitions change; freeze them or prove an
explicit bridge to the old semantics.
\item \textbf{Concrete representation preservation.}
For circuits retain input dimensions, fixed-leaf dimensions, register bases,
outputs, literal node order, topology, liveness, zero-default semantics,
and operation-count definitions. For witness maps retain coordinate order,
positivity, source binding, row permutations, and origin corrections.
\item \textbf{Strict arithmetic interfaces.}
Preserve natural-subtraction side conditions, strict digit/carry bounds,
and the exceptional all-maximum cyclic representation. Do not change
\code{word < radix^n - 1} into a non-strict inequality, or assume local
no-carry when the proof establishes a quotient fixed point instead.
\item \textbf{Independent specifications.}
The certified data must remain linked to the actual source polynomial,
literal DAG, and high-level recipe. A data hash proves identity, not
correctness; a table defined from itself proves neither independent
transcription nor specification agreement.
\end{enumerate}

Several useful simplifications are already present. The generic cyclic-carry
API and the periodicity/mod-four argument in
\src{NativeFieldCycles2011.lean}{49}{95} should be reused, not presented as
new refactors. Likewise the existing source-window, DAG soundness, power,
and degree interfaces are assets of the baseline, not deliverables of this
review.

\subsection{Reproduction commands and their limits}

The source-only checks below can be rerun from the extracted archive root:

\begin{Lean}
python3 verify_evidence.py
python3 validation/test_tools.py
python3 provenance/verify_sources.py
\end{Lean}

The upstream baseline commands, to run later in a configured full checkout,
include:

\begin{Lean}
lake build DiophantineConsistency
lake env lean Computability/DiophantineConsistency/Lean/checks/ZFCAxioms.lean
lake env lean Computability/DiophantineConsistency/Lean/checks/CompilerAxioms.lean
\end{Lean}

The supplied future replay harness takes a full clean pinned checkout,
an existing direct toolchain binary directory, and optionally a new output
directory. It installs nothing and keeps candidate outputs outside the
checkout:

\begin{Lean}
bash validation/replay_candidates.sh \
  /path/to/full/ProveIt /path/to/lean-4.32.0/bin /path/to/new-output
\end{Lean}

It targets Linux/WSL with GNU utilities. Even if it later succeeds, it only
establishes candidate elaboration and selected axiom checks, not a complete
downstream integration build or full regression pass. Its prerequisites and
the complete acceptance contract are in \pathcode{validation/README.md}.

When measuring performance, hold the toolchain, dependencies, machine,
parallelism, and cache state fixed. Record clean and incremental wall time,
peak resident memory, and generated bytes separately. An \code{.olean}
file's size is not proof-term size. No runtime, heartbeat, memory,
kernel-proof-size, or build-time saving was measured in this review.
